\def\stat{gon+zats}





\def\tit{ИНФОРМАЦИОННЫЕ ТРАНСФОРМАЦИИ ПАРАЛЛЕЛЬНЫХ ТЕКСТОВ В~ЗАДАЧАХ 


ИЗВЛЕЧЕНИЯ ЗНАНИЙ$^*$}





\def\titkol{Информационные трансформации параллельных текстов в~задачах 


извлечения знаний}








\def\aut{А.\,А.~Гончаров$^1$, И.\,М.~Зацман$^2$}





\def\autkol{А.\,А.~Гончаров, И.\,М.~Зацман}





\titel{\tit}{\aut}{\autkol}{\titkol}





\index{А.\,А.~Гончаров$^1$, И.\,М.~Зацман$^2$}


\index{Adamovich I.\,M.}


\index{Volkov O.\,I.}








{\renewcommand{\thefootnote}{\fnsymbol{footnote}}


\footnotetext[1]


{Работа выполнена в~Институте проблем информатики ФИЦ ИУ РАН при частичной поддержке РФФИ (проект 


18-07-00192).}}





\renewcommand{\thefootnote}{\arabic{footnote}}


\footnotetext[1]{Институт проблем информатики Федерального исследовательского центра <<Информатика 


и~управление>> Российской академии наук, \mbox{a.gonch48@gmail.com}}


\footnotetext[2]{Институт проблем информатики Федерального исследовательского центра <<Информатика 


и~управление>> Российской академии наук, \mbox{izatsman@yandex.ru}}








 \label{st\stat}





                  \thispagestyle{headings}


                  


  \vspace*{-12pt}


  


  


   \Abst{Рассматривается задача целенаправленного обнаружения и~заполнения лакун 


в~лингвистических типологиях, выступающих как формы представления знаний. Процесс 


решения этой задачи включает несколько многократно повторяемых стадий, которые 


в~совокупности образуют одну итерацию решения задачи целенаправленного извлечения 


знаний из параллельных текстов для заполнения лакун. Параллельные тексты как 


информационный ресурс трансформируются в~процессе решения этой задачи. Цель статьи 


состоит в~описании видов информационных трансформаций параллельных текстов, 


используемых в~начале процесса извлечения знаний и~заполнения лакун в~лингвистических 


типологиях, а~именно: фрагментации текстов на объекты интерпретации и~поиска тех из них, 


которые и~являются потенциальными источниками знания для заполнения лакун. В~статье 


процесс фрагментации позиционируется как один из видов информационных трансформаций 


параллельных текстов.}


   


   \KW{обнаружение лакун; заполнение лакун; лингвистические типологии; извлечение 


знаний из параллельных текстов; корпусная лингвистика; объекты интерпретации}








\DOI{10.14357\!/\!08696527190115} 








\vskip 10pt plus 9pt minus 6pt








      \thispagestyle{headings}


      


      \vspace*{-18pt} 





\section{Введение}





\vspace*{-2pt}





  Обнаружение и~заполнение лакун в~типологиях, рассматриваемых как формы 


представления лингвистических знаний,~--- это одна из задач проекта по гранту 


РФФИ, который в~настоящее время выполняется в~Институте проблем 


информатики ФИЦ ИУ РАН. Для решения данной задачи предлагается 


использовать лингвистическое аннотирование параллельных текстов~[1--4], 


которое вклю\-ча\-ет~5~мно\-го\-крат\-но повторяемых стадий. В~совокупности они 


образуют одну итерацию решения задачи целенаправленного извлечения 


знаний из параллельных текстов для заполнения лакун.


  


  В вышеупомянутом проекте предлагаемый подход к~извлечению знаний 


иллюстрируется на примере уточнения типологии значений немецких 


модальных глаголов и~их переводных эквивалентов на русском языке~[5]. Для 


проверки реализуемости разработанных методов обнаружения и~заполнения 


лакун используются тексты, полученные из не\-мец\-ко-рус\-ско\-го подкорпуса 


Национального корпуса русского языка (НКРЯ) общим объемом 2,6~млн 


словоупотреблений, включая 1,4~млн словоупотреблений в~оригинальных 


текстах на немецком языке и~1,2 млн словоупотреблений в~их переводах на 


русский язык.


  


  В работе~[6] дано описание вышеупомянутых 5~стадий обнаружения 


и~заполнения лакун:


  \begin{enumerate}[(1)]


\item определение одной или нескольких языковых единиц (ЯЕ), исследуемых 


на текущей итерации, а также поиск предложений, которые содержат 


ис\-сле\-ду\-емые ЯЕ, и~соответствующих им переводов в~{параллельных 


текстах};


  \item  лингвистическое аннотирование предложений с~исследуемыми ЯЕ и~их 


переводов, найденных в~параллельных текстах на первой стадии~[7];


  \item семантический анализ незавершенных аннотаций, помеченных 


рубриками, которые указывают на причины незавершенности аннотирования 


(в~том числе на лакуны в~типологии исследуемых ЯЕ);


  \item доработка аннотаций, включающих рубрики, которые указывают на 


лакуны в~типологии как на причину незавершенности, и~пополнение 


существующей типологии;


  \item вычисление значений числовых параметров, характеризующих 


состояние процесса пополнения типологии на момент завершения каждой 


итерации.


  \end{enumerate}


  


  Предложения и~их переводы, найденные на первой стадии, будем называть 


\textit{объектами интерпретации}. Три примера объектов интерпретации, 


состоящие из одного-двух предложений и~их переводов, приведены в~виде 


строк табл.~1.


  


\begin{table}\small


\begin{center}


\Caption{Примеры объектов интерпретации}


\vspace*{2ex}





\begin{tabular}{|p{63mm}|p{59mm}|}


\hline


--- ``Soll ich Ihr helfen, Jungfer?'' sagte ich.~--- Sie ward rot $\ddot{\mbox{u}}$ber und 


$\ddot{\mbox{u}}$ber.\newline


[Johann Wolfgang Goethe. Die Leiden des jungen Werther (1774)]&--- Помочь вам, девушка?~--- 


спросил~я. Она вся так и~зарделась.\newline


[Иоганн Вольфганг Гёте. Страдания юного Вертера (Н.~Касаткина,1954)]\\


\hline


Aber wie sollte ich von hier den Weg finden, der zur Menschenliebe 


f$\ddot{\mbox{u}}$hrte?\newline


[Hermann Hesse. Peter Camenzind (1904)]&Однако где же мне следовало искать путь, 


ведущий к~человеколюбию?\newline


[Герман Гессе. Петер Каменцинд (Р.~Эйвадис, 1995)]\\


\hline


--- ``Es kann sie jeder machen,'' versetzte sie drauf, ``und sollte denn in der weiten Welt kein 


M$\ddot{\mbox{a}}$dchen sein, das die W$\ddot{\mbox{u}}$nsche Ihres Herzens 


erf$\ddot{\mbox{u}}$llte?''\newline


[Johann Wolfgang Goethe. Die Leiden des jungen Werther (1774)]&---~Так всякий бы 


рассудил,~--- ответила она. --- Неужто во всём мире не найдётся девушки вам по 


сердцу?\newline


\newline


[Иоганн Вольфганг Гёте. Страдания юного Вертера (Н.~Касаткина, 1954)]\\


\hline


\end{tabular}


\end{center}


%\vspace*{-12pt}


\end{table}





  Эти примеры были найдены по лемме \textit{sollen}, так как до 


начала поиска было решено на текущей итерации исследовать конструкции 


с~немецким модальным глаголом \textit{sollen} и~их переводы на русский 


язык~[5]. Ввиду того что данный глагол обладает ярко выраженной полисемией 


(в~\cite{8-gz} для него описано~13~значений), для уточнения предлагаемой 


типологии значений были использованы немецко-русские параллельные тексты 


преимущественно художественных произведений, в~которых всего было 


найдено более двух тысяч объектов интерпретации, содержащих словоформы 


леммы \textit{sollen}~\cite{6-gz}. %Три из них приведены в~табл.~1.


  


  Используемые в~проекте параллельные тексты выровнены (оригинальным 


предложениям поставлены в~соответствие их переводы, см.\ строки табл.~1), т.\,е.\


 в~процессе формирования НКРЯ была выполнена их фрагментация на 


строки с~предложениями на немецком и~русском языке. С~точки зрения 


информационных трансформаций можно говорить о~том, что до начала 


итерационного выполнения~5~перечисленных стадий была осуществлена 


фрагментация оригинальных текстов и~соответствующих им переводов на 


объекты интерпретации.


  


  Однако после выполнения поиска появляется возможность переформировать 


любой найденный объект интерпретации путем добавления к~нему одной или 


нескольких строк, расположенных соответственно в~тексте и~переводе 


произведения до этого объекта и\!/\!или после него. Например, если для 


интерпретации выбрана некоторая найденная строка, то она может быть 


объединена с~соседними строками (необязательно смежными по тексту книги) 


в~новый объект интерпретации, что обеспечивается за счет расширения границ 


контекста, необходимого для выполнения аннотирования.


  


  Цель статьи состоит в~описании информационных трансформаций 


параллельных текстов, обеспечивающих поиск и~переформирование именно тех 


объектов интерпретации, которые являются потенциальными источниками 


знания для уточнения типологии в~процессе обнаружения и~заполнения лакун.


  


\section{Формирование и~поиск объектов интерпретации}





  Формирование объектов интерпретации с~исследуемыми ЯЕ выполняется 


в~три этапа:


  \begin{enumerate}[(1)]


  \item  выравнивание параллельных текстов;


  \item поиск объектов интерпретации с~исследуемыми ЯЕ, которые являются 


потенциальными источниками знания для уточнения типологии (первая стадия 


итерации);


  \item переформирование найденных объектов интерпретации (посредством 


расширения границ объекта интерпретации за счет других строк таблицы, 


содержащих текст и~перевод произведения), если для аннотирования требуется 


более широкий контекст (вторая стадия итерации).


  \end{enumerate}


  


  Результатом первого этапа формирования объектов интерпретации являются 


пары фрагментов текста, включающие предложения оригинала 


и~соответствующие им переводы. В~рамках надкорпусной базы данных (НБД) на 


текущий момент используются полные тексты произведений, которые 


визуально представлены в~форме таблицы. Каждая строка таблицы содержит 


фрагмент текста оригинала (в~левом столбце) и~соответствующий ему 


фрагмент текста перевода (в~правом столбце), фрагменты же в~большинстве 


случаев состоят из одного предложения~\cite{9-gz, 10-gz}. Если в~табл.~1 


убрать названия произведений, имена авторов и~переводчиков, то ее строки 


будут примерами результатов выполнения первого этапа.


  


  Кратко опишем технологию выравнивания, исходя из того, что тексты 


оригинала и~перевода уже представлены в~электронном варианте. Сначала 


тексты проверяются на предмет опечаток, что можно сделать, например, при 


помощи функции <<Правописание>> MS~Word. Затем в~текстах удаляется 


нумерация страниц, знаки табуляции заменяются пробелами, разрывы 


страниц~--- знаками абзаца, лишние пробелы и~мягкие переносы удаляются, 


знаки тире унифицируются. Цель этих операций~--- нормализация знаков 


форматирования, что в~дальнейшем поможет избежать сбоев при 


использовании программы выравнивания параллельных текстов. 


Необходимо также учитывать особенности конкретных языков. Так,


в~текстах на французском языке нередко приходится заменять сочетание букв 


<<oe>> на <<c\!e>>.


  


  При подготовке текстов к~выравниванию встает вопрос о необходимости 


сохранения и~редактирования постраничных сносок. Если принимается 


решение об их сохранении, то текст сноски заключается в~фигурные скобки 


и~помещается после предложения, к~которому он относится. При этом 


неинформативные сноски, как правило, удаляются.


  


  После всех вышеперечисленных операций выполняется выравнивание 


предложений оригинального текста и~его перевода при помощи программы 


<<Евклид>>~\cite{11-gz}, созданной на основе находящейся в~открытом 


доступе программы выравнивания текстов HunAlign~\cite{12-gz}. Программа 


<<Евклид>> позволяет установить соответствие между одним или несколькими 


предложениями текста оригинала и~их переводами. В~автоматическом режиме 


она формирует таблицу с~полным текстом произведения в~левом столбце и~его 


переводом в~правом.


  


  Полученная таблица может содержать ошибки выравнивания, которые 


исправляются вручную при помощи нескольких команд. Во-пер\-вых, это 


перенос фрагмента текста в~таблице на строку выше или ниже. Во-вто\-рых, это 


добавление пустой строки в~таблицу или, напротив, удаление лишней (пустой) 


строки. В-треть\-их, это распределение фрагментов текста оригинала 


и~перевода, помещенных в~ходе автоматического выравнивания в~одну строку, 


по двум или более новым строкам, а~также объединение  


фрагментов текста из нескольких строк. Следует отметить, что программа <<Евклид>> 


в~автоматическом режиме выделяет красным цветом те фрагменты, где 


с~большой вероятностью могут присутствовать ошибки выравнивания. Таким 


образом, с~помощью этой программы линейный текст и~его перевод 


трансформируются в~таблицу, состоящую из двух столбцов. Предложения 


текста оригинала распределяются по строкам в~левом столбце этой таблицы (по 


одному или более предложений в~строке), а~перевод этих предложений~--- в~правом.


  


  После информационной трансформации текста произведения и~его перевода 


в~таблицу выполняется разметка выровненных текстов программами 


Яндекса~\cite{13-gz, 14-gz} согласно Морфологическому стандарту НКРЯ. Перечень 


тегов, используемых при разметке, приведен на странице по адресу: {\small{\sf 


http://www. ruscorpora.ru/corpora-morph.html}}. Морфологическая разметка 


представляет собой информационную трансформацию, которая преобразует 


таблицу выровненных текстов в~таблицу морфологически размеченных текстов 


произведения и~его перевода. Например, три первые слова оригинального 


текста первого объекта интерпретации в~табл.~1 (Soll ich Ihr) будут размечены 


следующим образом:





\medskip





\noindent


 {\footnotesize


%  \begin{verbatim}


   \ttfamily


<w>\\


<ana lex="\underline{soll}" sem="" disamb="yes" gr="S,n,nom,norm,sg" sem2=""/>\\


<ana lex="\underline{soll}" sem="" disamb="yes" gr="S,dat,n,norm,sg" sem2=""/>\\


<ana lex="\underline{soll}" sem="" disamb="yes" gr="S,acc,n,norm,sg" sem2=""/>\\


<ana lex="\underline{sollen}" sem="" gr="V,1p,act,indic,intr,norm,praes,sg" sem2=""/>\\


<ana lex="\underline{sollen}" sem="" gr="V,3p,act,indic,intr,norm,praes,sg"\\


\hspace*{10mm}sem2=""/>{\bf Soll}</w>\\


<w>\\


<ana lex="\underline{ich}" sem="" disamb="yes" gr="SPRO,1p,nom,norm,sg" sem2=""/>{\bf ich}</w>\\


<w>\\


<ana lex="\underline{ihr}" sem="" disamb="yes" gr="SPRO,2p,nom,norm,pl" sem2=""/>\\


<ana lex="\underline{sie}" sem="" gr="SPRO,3p,dat,f,norm,sg" sem2=""/>\\


<ana lex="\underline{ihr}" sem="" gr="APRO,m,nom,norm,sg" sem2=""/>\\


<ana lex="\underline{ihr}" sem="" gr="APRO,n,nom,norm,sg" sem2=""/>\\


<ana lex="\underline{ihr}" sem="" gr="APRO,acc,n,norm,sg" sem2=""/>{\bf Ihr}</w> 


%\end{verbatim}


}


  


  \medskip


  


 \noindent


  Здесь приведен фрагмент морфологически размеченного текста с~неснятой 


грамматической омонимией. Это означает, что для каждой словоформы 


пе\-ре\-чис\-ле\-ны все возможные варианты интерпретации. Видимые пользователю 


единицы для наглядности выделены в~примере полужирным шрифтом, 


а~леммы, которые могут соответствовать каждой из этих трех словоформ,~--- 


подчеркиванием; кроме того, сплошной текст для удобства разбит на строки.


  


  Так, для словоформы \textit{soll} представлено 5~вариантов 


интерпретации. 


  


  Три из них соответствуют лемме \textit{Soll} с~разными наборами 


морфологических характеристик:


  \begin{enumerate}[(1)]


\item существительное (S), средний род (n), именительный падеж (nom), 


отсутствие особенностей словоформы (norm), единственное число (sg);


\item существительное (S), дательный падеж (dat), средний род (n), отсутствие 


особенностей словоформы (norm), единственное число (sg);


\item существительное (S), винительный падеж (acc), средний род (n), 


отсутствие особенностей словоформы (norm), единственное число (sg).


\end{enumerate}





  Два других варианта разметки включают морфологические характеристики 


омонимичных форм модального глагола \textit{sollen}:


  \begin{enumerate}[(1)]


  \setcounter{enumi}{3}


\item глагол (V), первое лицо (1p), действительный залог (act), изъявительное 


наклонение (indic), непереходность (intr), отсутствие особенностей словоформы 


(norm), настоящее время (praes), единственное число (sg);


\item глагол (V), третье лицо (3p), действительный залог (act), изъявительное 


наклонение (indic), непереходность (intr), отсутствие особенностей словоформы 


(norm), настоящее время (praes), единственное число (sg).


\end{enumerate}





  В корпусе со снятой омонимией для данного словоупотребления был бы 


сохранен лишь вариант разметки под номером~4.


  


  Для словоформы \textit{ich} вариант разметки лишь один (ввиду отсутствия 


омонимичных словоформ): местоимение-существительное (SPRO), первое лицо 


(1p), именительный падеж (nom), отсутствие особенностей словоформы (norm), 


единственное число (sg).


  


  Для словоформы \textit{Ihr} представлены 5 вариантов интерпретации для 


трех лемм.


  


  Один вариант для леммы \textit{ihr}:


  \begin{enumerate}[(1)]


\item местоимение-существительное (SPRO), второе лицо (2p), именительный 


падеж (nom), отсутствие особенностей словоформы (norm), множественное 


число (pl).


\end{enumerate}





  Один вариант для леммы \textit{sie}:


  \begin{enumerate}[(1)]


  \setcounter{enumi}{1}


\item местоимение-существительное (SPRO), третье лицо (3p), дательный падеж 


(dat), женский род (f), отсутствие особенностей словоформы (norm), 


единственное число (sg).


\end{enumerate}





  Три варианта для леммы \textit{ihr}, омонимичной первой:


  \begin{enumerate}[(1)]


    \setcounter{enumi}{2}


\item местоимение-прилагательное (APRO), мужской род (m), именительный 


падеж (nom), отсутствие особенностей словоформы (norm), единственное число 


(sg);


\item местоимение-прилагательное (APRO), средний род (n), 


именительный падеж (nom), отсутствие особенностей словоформы 


(norm), единственное число (sg);


\item местоимение-прилагательное (APRO), винительный падеж (acc), средний 


род (n), отсутствие особенностей словоформы (norm), единственное число (sg).


\end{enumerate}





  В корпусе со снятой омонимией для данного словоупотребления был бы 


сохранен лишь вариант разметки под номером~2.


  


  Аналогичным образом размечается и~перевод на русский язык. При этом 


сохраняется выравнивание, полученное в~результате предыдущей 


информационной трансформации. Морфологическая разметка текстов дает 


возможность осуществлять поиск не только по точной форме, но также по 


лемме или по грамматическим характеристикам.


  


  Морфологически размеченные тексты загружаются в~НБД, 


  функциональность которой, включая поисковые возможности, 


описана на примерах нескольких категорий ЯЕ  


в~работах~\cite{15-gz, 16-gz, 17-gz, 18-gz, 19-gz, 20-gz, 21-gz, 22-gz}. Приведем 


пример поиска объектов интерпретации в~НБД, в~которую были загружены 


выровненные немецко-русские параллельные тексты общим объемом 2,6~млн 


словоупотреблений. Таблица загруженных текстов состоит из 83\,190~строк, 


включая более двух тысяч объектов интерпретации~\cite{6-gz}, содержащих 


словоформы леммы \textit{sollen}, на примере которой иллюстрируются 


поисковые возможности НБД.


  


  Для извлечения объектов интерпретации в~НБД реализован поиск строк 


в~таблице как по словоформам, так и~по леммам. Поиск можно вести как по 


тексту одного произведения, так и~по всем загруженным текстам. Например, 


если задать поиск строк во всех произведениях по лемме \textit{sollen}, то будет 


найдена~2041~строка с~исследуемой ЯЕ, 4~из которых приведены 


в~табл.~2\footnote{Все четыре объекта интерпретации из книги 


\Au{М.~Энде}. Момо (перевод Ю.\,И.~Коринец, 1982).}.





\begin{table}\small %tabl2


\begin{center}


\Caption{Примеры результатов поиска пар фрагментов текста по лемме \textit{sollen}}


\vspace*{2ex}





\begin{tabular}{|c|p{53mm}|p{53mm}|}


\hline


\tabcolsep=0pt\begin{tabular}{c}Номер\\ строки\end{tabular}&\multicolumn{1}{c|}{Оригинальный текст}&\multicolumn{1}{c|}{Перевод}\\


\hline


1479&Oder meinst du, ich soll's?&Или ты думаешь~--- должен?\\


\hline


1582&Soll ich es dir zeigen?&Показать тебе?\\


\hline


1636&Du sollst nur damit spielen, so wie ich es dir erkl$\ddot{\mbox{a}}$rt habe.&Только 


играть должна ты со всем этим так, как я тебе объяснил.\\


\hline


1660&``Ich glaube,'' erwiderte er sanft, ``wir sollten einmal ernsthaft miteinander reden, 


Kleine, damit du lernst, worauf es ankommt.''&---~Я~думаю, что нам надо серьёзно 


поговорить,~--- сказал он.~--- Ты должна понять, что к~чему.\\


\hline


\end{tabular}


\end{center}


\vspace*{-12pt}


\end{table}


  


  Найденные строки по умолчанию являются объектами интерпретации, 


содержащими исследуемые ЯЕ и~представляющими собой потенциальные 


источники знания для уточнения типологии.


  


\section{Расширение границ объектов интерпретации}





  Лингвист, выполняющий аннотирование, может при необходимости 


переформировать найденный объект интерпретации, объединяя строки 


таблицы. Примером может послужить строка №\,1479 (см.\ табл.~2). 


Содержащийся в~ней фрагмент текста недостаточен для формирования 


аннотации, так как он не позволяет описать функционирование модального 


глагола \textit{sollen} (словоформа \textit{soll}) и~определить его семантику. 


Поэтому в~данном случае для обоснования выбора одного из словарных 


значений исследуемой ЯЕ в~процессе аннотирования необходимо расширить 


контекст за счет строк 1476--1478, с~которыми строка~1479 тесно связана по 


смыслу, и~таким образом сформировать новый объект интерпретации 


(включающий все строки табл.~3). Это даст возможность определить 


релевантное значение глагола \textit{sollen}~--- обязанность что-либо делать по 


чьему-либо указанию, по закону, по правилам и~т.\,п.~\cite{8-gz}.


  


  Таким образом, при недостаточности контекста в~процессе аннотирования 


может возникать потребность в~его расширении, в~первую очередь для 


определения значения исследуемой ЯЕ и~особенностей ее функционирования. 


В~этих случаях возможность такой информационной трансформации 


выровненных и~размеченных текстов оригинала и~перевода в~НБД, как 


расширение границ объектов интерпретации, сложно переоценить. Эта функция 


широко используется в~процессе лингвистического аннотирования не только 


модальных глаголов, но и~других категорий ЯЕ, например коннекторов, в~сферу 


действия которых может входить не только часть предложения или отдельное 


предложение, но и~фрагмент текста, включающий два и~более 


предложений~\cite[с.~17--20]{23-gz}.


  


\begin{table}\small %tabl3


\begin{center}


\Caption{Объединение соседних строк в~один объект интерпретации}


\vspace*{2ex}





\tabcolsep=2.5pt


\begin{tabular}{|c|l|l|}


\hline


1476&Ich \mbox{wei{\!\!\ptb{\ss}}} wirklich nicht, was ich tun soll.&Я правда не знаю, что 


делать.\\


1477&Aber alle machen's doch heute so.&Но сегодня все так поступают.\\


1478&Warum soll ich allein es anders machen?&Почему я должен поступать по-иному?\\


1479&\textbf{Oder meinst du, ich soll's?}&\textbf{Или ты думаешь~--- должен?}\\


\hline


\end{tabular}


\end{center}


%\vspace*{-12pt}


\end{table}





  Отметим, что в~описанном выше расширении контекста для строки~1479 


в~расширенный объект интерпретации попали еще два случая употребления 


исследуемой ЯЕ (словоформа \textit{soll} в~строках~1476 и~1478). И~если 


контекст строки~1476 достаточен для определения значения глагола 


\textit{sollen} (выражение неуверенности, растерянности в~главных 


и~придаточных вопросительных предложениях с~вопросительными 


словами~\cite{8-gz}), то для аннотирования исследуемой ЯЕ в~строке~1478 


снова возникает необходимость расширения контекста. Расширенный объект 


интерпретации будет включать строки~1476--1478. Таким образом, этот объект 


интерпретации будет вложен в~объект интерпретации, включающий 


строки~1476--1479. Однако в~общем случае расширенные объекты 


интерпретации могут пересекаться по составляющим их строкам и~тогда 


границы таких объектов интерпретации нельзя описать, используя 


расширяемый язык разметки (eXtensible Markup Language~--- 


XML)\footnote{Согласно синтаксическим правилам XML любой 


элемент, начинающийся внутри другого элемента, должен 


заканчиваться внутри элемента, в~котором он начался.}.


 Для 


информационной трансформации выровненных и~размеченных текстов 


оригинала и~перевода, состоящей в~расширении границ объектов 


интерпретации, используются функции НБД, описанные в~работе~\cite{24-gz}.





\section{Заключение}





  Подход к~извлечению знаний из параллельных текстов, используемый 


в~проекте, предполагает, что сначала выполняется ряд информационных 


трансформаций этих текстов, включая фрагментацию на объекты 


интерпретации и~поиск тех из них, которые и~являются потенциальными 


источниками знания для формирования и~уточнения лингвистических 


типологий. В~статье рассмотрены следующие информационные 


трансформации параллельных текстов:


  \begin{enumerate}[(1)]


  \item преобразование линейного текста произведения и~его перевода 


в~таблицу, состоящую из двух столбцов. Предложения текста произведения 


распределяются по строкам в~левом столбце этой таблицы (по одному или 


более предложений в~строке). Перевод предложений из левого столбца строки 


размещается в~правом столбце той же строки. В~результате такой 


трансформации соответствие оригинальных предложений и~их переводов, 


основанное на эквивалентности смыслового содержания, фиксируется 


в~табличной \mbox{форме};


  \item преобразование таблицы выровненных текстов в~таблицу текстов, 


размеченных согласно Морфологическому стандарту НКРЯ. При этом 


сохраняется выравнивание, полученное в~результате первой информационной 


трансформации. Разметка, включенная в~тексты, дает возможность искать 


объекты интерпретации не только в~точной форме, но также по лемме и~по 


грамматическим характеристикам (в~оригинальном тексте и\!/\!или в~его 


переводе);


  \item переформирование объектов интерпретации, являющихся 


потенциальными источниками знания для уточнения лингвистических 


типологий в~процессе обнаружения и~заполнения лакун, которое 


обеспечивается за счет расширения границ контекста исследуемых ЯЕ, 


необходимого для выполнения аннотирования ЯЕ и~их контекстов.


  \end{enumerate}


  


  В~заключение отметим, что используемый в~проекте метод аннотирования 


объектов интерпретации, содержащих исследуемые ЯЕ, характеризуется тем, что 


допускает возможность незавершенности процесса аннотирования. Именно эта 


возможность и~позволяет фиксировать лакуны в~уточняемой типологии 


и~заполнять их в~процессе доработки незавершенных  


аннотаций~\cite{6-gz, 24-gz}.


     


{\small\frenchspacing


{%\baselineskip=10.8pt


\addcontentsline{toc}{section}{Литература}


\begin{thebibliography}{99}





\bibitem{3-gz} %1


\Au{Дурново А.\,А., Зацман И.\,М., Лощилова~Е.\,Ю.} Кросслингвистическая база данных для 


аннотирования логико-семантических отношений в~тексте~/\!\!/ Системы и~средства 


информатики, 2016. Т.~26. №\,4. С.~124--137.


\bibitem{4-gz} %2


\Au{Зацман И.\,М., Мамонова О.\,С., Щурова~А.\,Ю.} Обратимость и~альтернативность 


генерализации моделей перевода коннекторов в~параллельных текстах~/\!\!/ Системы 


и~средства информатики, 2017. Т.~27. №\,2. С.~125--142.





\bibitem{2-gz} %3


\Au{Зацман И.\,М., Кружков~М.\,Г., Лощилова~Е.\,Ю.} Методы анализа частотности моделей 


перевода коннекторов и~обратимость генерализации статистических данных~/\!\!/ Системы 


и~средства информатики, 2017. Т.~27. №\,4. С.~164--176.





\bibitem{1-gz} %4


\Au{Kruzhkov M.\,G.} Approaches to annotation of discourse relations in linguistic corpora~/\!\!/ 


Информатика и~её применения, 2017. Т.~11. Вып.~4. С.~118--125.





\bibitem{5-gz}


\Au{Добровольский Д.\,О., Зализняк~Анна~А.} Немецкие конструкции с~модальными 


глаголами и~их русские соответствия: проект надкорпусной базы данных~/\!\!/ 


Компьютерная лингвистика и~интеллектуальные технологии: По мат-лам Междунар. конф. 


<<Диалог>>.~--- М.: РГГУ, 2018. Вып.~17(24). С.~172--184.


\bibitem{6-gz}


\Au{Зацман И.\,М.} Стадии целенаправленного извлечения знаний, имплицированных 


в~параллельных текстах~/\!\!/ Системы и~средства информатики, 2018. Т.~28. №\,3.  


С.~175--188.


\bibitem{7-gz}


Handbook of linguistic annotation~/ Eds. N.~Ide, J.~Pustejovsky.~--- Dordrecht, The Netherlands: 


Springer Science\;+\;Business Media, 2017. 1468~p.


\bibitem{8-gz}


Немецко-русский словарь: актуальная лексика~/ Под ред. Д.\,О.~Добровольского.~--- М.: 


Лексрус, 2019 (в печати).


\bibitem{9-gz}


\Au{Добровольский Д.\,О., Кретов~А.\,А., Шаров~С.\,А.} Корпус параллельных текстов~/\!\!/ 


На\-уч\-но-тех\-ни\-че\-ская информация. Серия~2: Информационные процессы и~системы, 


2005. №\,6. С.~16--27.


\bibitem{10-gz}


\Au{Loiseau S., Sitchinava~D.\,V., Zalizniak~Anna~A., Zatsman~I.\,M.} Information technologies 


for creating the database of equivalent verbal forms in the Russian-French multivariant parallel 


corpus~/\!\!/ Информатика и~её применения, 2013. Т.~7. Вып.~2. С.~100--109.


\bibitem{11-gz}


\Au{Сичинава Д.\,В.} Параллельные тексты в~составе Национального корпуса русского языка: 


новые направления развития и~результаты~/\!\!/ Труды Института русского языка им.\ 


В.\,В.~Виноградова, 2015. №\,6. С.~194--235.


\bibitem{12-gz}


\Au{Varga D., N$\acute{\mbox{e}}$meth~L., Hal$\acute{\mbox{a}}$csy~P., Kornai~A., 


Tr$\acute{\mbox{o}}$n~V., Nagy~V.} Parallel corpora for medium density languages~/\!\!/  


Conference (International) on Recent Advances in Natural Language Processing  Proceedings.~--- 


Shoumen, Bulgaria: INCOMA Ltd., 2005. P.~590--596.


\bibitem{13-gz}


\Au{Segalovich I.} A~fast morphological algorithm with unknown word guessing induced by 


a~dictionary for a~Web Search Engine~/\!\!/ Conference (International) on Machine Learning; 


Models, Technologies and Applications Proceedings.~--- Las Vegas, NV, USA: CSREA Press, 2003.  


P.~273--280.


\bibitem{14-gz}


\Au{Зобнин А.\,И., Носырев Г.\,В.} Морфологический анализатор MyStem~3.0~/\!\!/ Труды 


Института русского языка им.\ В.\,В.~Виноградова, 2015. №\,6. С.~300--310.





\bibitem{22-gz} %15


\Au{Zalizniak Anna~A., Sitchinava~D.\,V., Loiseau~S., Kruzhkov~M., Zatsman~I.\,M.} Database of 


equivalent verbal forms in a Russian-French multivariant parallel corpus~/\!\!/ 


Conference 


(International) on Artificial Intelligence Proceedings.~---  Las Vegas,


NV, USA: CSREA Press, 2013. Vol.~1. 


P.~101--107.











\bibitem{20-gz} %16


\Au{Kruzhkov M.\,G., Buntman~N.\,V., Loshchilova~E.\,Ju., Sitchinava~D.\,V., Zalizniak~Anna~A., 


Zatsman~I.\,M.} A~database of Russian verbal forms and their French translation equivalents~/\!\!/ 


Компьютерная лингвистика и~интеллектуальные технологии: по мат-лам ежегодной 


Междунар. конф. <<Диалог>>.~--- М.: РГГУ, 2014. Вып.~13(20). С.~275--287.


\bibitem{21-gz} %17


\Au{Бунтман~Н.\,В., Зализняк~Анна~А., Зацман~И.\,М., Кружков~М.\,Г., Лощилова~Е.\,Ю., 


Сичинава~Д.\,В.} Информационные технологии корпусных исследований: принципы 


построения кросслингвистических баз данных~/\!\!/ Информатика и~её применения, 2014. 


Т.~8. Вып.~2. С.~98--110.





\bibitem{17-gz} %18


\Au{Попкова Н.\,А., Инькова~О.\,Ю., Зацман~И.\,М., Кружков~М.\,Г.} Методика построения 


моноэквиваленций в~надкорпусной базе данных коннекторов~/\!\!/ Задачи современной 


информатики: Тр. Второй молодежной научн. конф.~--- М.: ФИЦ ИУ РАН, 2015. С.~143--153.


\bibitem{18-gz} %19


\Au{Зализняк Анна А., Зацман~И.\,М., Инькова~О.\,Ю., Кружков~М.\,Г.} Надкорпусные базы 


данных как лингвистический ресурс~/\!\!/ Корпусная лингвистика-2015: Тр. Междунар. 


конф.~--- СПб.: СПбГУ, 2015. С.~211--218.


\bibitem{19-gz} %20


\Au{Кружков М.\,Г.} Информационные ресурсы контрастивных лингвистических 


исследований: электронные корпуса текстов~/\!\!/ Системы и~средства информатики, 2015. 


Т.~25. №\,2. С.~140--159.





\bibitem{15-gz} %21


\Au{Инькова О.\,Ю., Кружков~М.\,Г.} Надкорпусные рус\-ско-фран\-цуз\-ские базы данных 


глагольных форм и~коннекторов~/\!\!/ Lingue slave a~confronto~/ Eds. O.~Inkova,  


A.~Trovesi.~--- Bergamo: Bergamo University Press, 2016. P.~365--392.





\bibitem{16-gz} %22


\Au{Зализняк Анна А., Кружков~М.\,Г.} База данных безличных глагольных конструкций 


русского языка~/\!\!/ Информатика и~её применения, 2016. Т.~10. Вып.~4. С.~132--141.











\bibitem{23-gz}


\Au{Инькова-Манзотти О.\,Ю.} Коннекторы противопоставления во французском и~русском 


языках (сопоставительное исследование).~--- М.: Информэлектро, 2001. 434~с.


\bibitem{24-gz}


\Au{Зацман И.\,М.} Целенаправленное развитие систем лингвистических знаний: выявление 


и~заполнение лакун~/\!\!/ Информатика и~её применения, 2019. Т.~13. Вып.~1. С.~91--98.


        \end{thebibliography}


} }








\vspace*{-3pt}





\hfill{\small\textit{Поступила в~редакцию 15.02.19}}














\vspace*{8pt}





\hrule





\vspace*{2pt}





%\newpage





%\vspace*{-24pt}





\hrule





%\vspace*{9pt}











\def\tit{INFORMATION TRANSFORMATIONS OF PARALLEL TEXTS  


IN~KNOWLEDGE EXTRACTION}





\def\titkol{Information transformations of parallel texts  


in~knowledge extraction}








\def\aut{A.\,A.~Goncharov and~I.\,M.~Zatsman}


\def\autkol{A.\,A.~Goncharov and~I.\,M.~Zatsman}











\titel{\tit}{\aut}{\autkol}{\titkol}





\vspace*{-9pt}





\noindent


   Institute of Informatics Problems, Federal Research Center ``Computer 


Science and Control'' of the Russian Academy of Sciences, 44-2~Vavilov Str., 


Moscow 119133, Russian Federation





 \noindent





\def\leftfootline{\small{\textbf{\thepage}


\hfill Sistemy i Sredstva Informatiki~---


Systems and Means of Informatics 2019 vol~29 no\,1}


}%


 \def\rightfootline{\small{Sistemy i Sredstva Informatiki~---


 Systems and Means of Informatics   2019   vol~29   no\,1


\hfill \textbf{\thepage}}}

















   \Abste{The paper examines the task of goal-oriented discovery and filling of lacunas in 


linguistic typologies considered as forms of knowledge representation. The process of solving this task 


includes several repeated stages which collectively form one iteration of the proposed solution to 


the task of goal-oriented knowledge discovery in parallel texts required to fill the lacunas. Parallel 


texts as an


information resource are transformed in the process of solving this task. The


purpose of 


the paper is to describe the types of information transformations of parallel texts that are used 


during early stages of the process of knowledge\linebreak\vspace*{-12pt}}





\Abstend{discovery and filling of lacunas in linguistic 


typologies. As a~part of knowledge discovery, first,


the parallel texts are fragmented into objects of 


interpretation and then, the search for potential sources of knowledge capable to fill the lacunas is 


performed. This paper considers this fragmentation process as one of the information transformation 


types of parallel texts.}


   


   \KWE{discovery of lacunas; filling of lacunas; linguistic typology; knowledge extraction from 


parallel texts; corpus linguistics; objects of interpretation}


  


 \DOI{10.14357\!/\!08696527190115}





%\vspace*{-24pt}





\Ack


\noindent


The work was performed at the Institute of 


Informatics Problems, Federal Research Center ``Computer 


Science and Control'' of the Russian Academy of Sciences


with partial financial support of the Russian Foundation for


Basic Research  (project  18-07-00192).





 








\renewcommand{\bibname}{\protect\rmfamily References}


%\renewcommand{\bibname}{\large\protect\rm References}





{\small\frenchspacing


{%\baselineskip=10.8pt


\addcontentsline{toc}{section}{References}


\begin{thebibliography}{99}


\bibitem{3-gz-1} %1


\Aue{Durnovo, A.\,A., I.\,M.~Zatsman, and E.\,Yu.~Loshchilova.} 2016.  


Krosslingvi\-sti\-che\-skaya baza dannykh dlya annotirovaniya logiko-semanticheskikh 


otnosheniy v~tekste [Cross-lingual database for annotating logical-semantic relations 


in the text]. \textit{Sistemy i~Sredstva Informatiki~--- Systems and Means of 


Informatics} 26(4):124--137.


\bibitem{4-gz-1} %2


\Aue{Zatsman, I.\,M., O.\,S. Mamonova, and A.\,Yu.~Shchurova.} 2017. Obratimost' 


i~al'ternativnost' generalizatsii modeley perevoda konnektorov v~parallel'nykh 


teks\-takh [Reversibility and alternativeness of generalization of connectives 


translations models in parallel texts]. \textit{Sistemy i~Sredstva Informatiki~--- 


Systems and Means of Informatics} 27(2):125--142.





\bibitem{2-gz-1} %3


\Aue{Zatsman, I.\,M., M.\,G. Kruzhkov, and E.\,Yu.~Loshchilova.} 2017. Metody 


analiza chastotnosti modeley perevoda konnektorov i~obratimost' generalizatsii 


statisticheskikh dannykh [Methods of frequency analysis of connectives translations 


and reversibility of statistical data generalization]. \textit{Sistemy i~Sredstva 


Informatiki~--- Systems and Means of Informatics} 27(4):164--176.





\bibitem{1-gz-1} %4


\Aue{Kruzhkov, M.\,G.} 2017. Approaches to annotation of discourse relations in 


linguistic corpora. \textit{Informatika i ee primeneniya~--- Inform. Appl.}  


11(4):118--125.





\bibitem{5-gz-1}


\Aue{Dobrovol'skiy, D.\,O., and Anna A.~Zaliznyak.} 2018. Nemetskie konstruktsii 


s~modal'nymi glagolami i~ikh russkie sootvetstviya: proekt nadkorpusnoy bazy 


dannykh [German constructions with modal verbs and their Russian correlates: 


A~supracorpora database project]. \textit{Komp'yuternaya lingvistika 


i~intellektual'nye tekhnologii. Po mat-lam ezhegodnoy Mezhdunar. konf. ``Dialog.''} 


[Computational Linguistics and Intellectual Technologies. Papers from the Annual 


Conference (International) ``Dialogue'']. Moscow: RGGU. 17(24):172--184.


\bibitem{6-gz-1}


\Aue{Zatsman, I.\,M.} 2018. Stadii tselenapravlennogo izvlecheniya znaniy, 


implitsirovannykh v~parallel'nykh tekstakh [Stages of goal-oriented discovery of 


knowledge implied in parallel texts]. \textit{Sistemy i~Sredstva Informatiki~--- 


Systems and Means of Informatics} 28(3):175--188.


\bibitem{7-gz-1}


Ide,~N., and J.~Pustejovsky, eds.  2017.


\textit{Handbook of linguistic annotation}. 


Dordrecht, The Netherlands: Springer Science\;+\;Business Media. 1468~p.


\bibitem{8-gz-1}


Dobrovol'skiy, D.\,O., ed. 2019 (in press). \textit{Nemetsko-russkiy slovar': 


aktual'naya lek\-si\-ka} [German-Russian dictionary: Actual vocabulary]. Moscow: 


Leksrus.


\bibitem{9-gz-1}


\Aue{Dobrovol'skiy, D.\,O., A.\,A.~Kretov, and S.\,A.~Sharov.} 2005. Korpus 


parallel'nykh teks\-tov [Corpus of parallel texts]. \textit{Automatic


Documentation Math. Linguistics}\linebreak 6:16--27.


\bibitem{10-gz-1}


\Aue{Loiseau, S., D.\,V.~Sitchinava, Anna A.~Zalizniak, and I.\,M.~Zatsman.} 


2013. Information technologies for creating the database of equivalent verbal forms 


in the Russian-French multivariant parallel corpus. \textit{Informatika i~ee 


Primeneniya~--- Inform. Appl.}7(2):100--109.


\bibitem{11-gz-1}


\Aue{Sitchinava, D.\,V.} 2015. Parallel'nye teksty v~sostave Natsional'nogo korpusa 


russkogo yazyka: novye napravleniya razvitiya i~rezul'taty [Parallel texts within the 


Russian National Corpus: New directions and results]. \textit{Trudy Instituta russkogo 


yazyka im.\ V.\,V.~Vinogradova} [Proceedings of the V.\,V.~Vinogradov Russian 


Language Institute]. 6:194--235.


\bibitem{12-gz-1}


\Aue{Varga D., L. N$\acute{\mbox{e}}$meth, P. Hal$\acute{\mbox{a}}$csy, 


A.~Kornai, V.~Tr$\acute{\mbox{o}}$n, and V.~Nagy.} 2005. Parallel corpora for 


medium density languages. \textit{Conference (International) 


on Recent Advances in Natural Language Processing  Proceedings}. 


Shoumen, Bulgaria: INCOMA Ltd. 590--596.


\bibitem{13-gz-1}


\Aue{Segalovich, I.} 2003. A~fast morphological algorithm with unknown word 


guessing induced by a~dictionary for a~Web Search Engine. \textit{Conference 


(International) on Machine Learning: Models, Technologies and Applications 


Proceedings}. Las Vegas, NV: CSREA Press. 273--280.


\bibitem{14-gz-1}


\Aue{Zobnin, A.\,I., and G.\,V.~Nosyrev.} 2015. Morfologicheskiy analizator MyStem~3.0 


[Morphological analyzer MyStem~3.0]. \textit{Trudy Instituta russkogo yazyka im.\ 


V.\,V.~Vinogradova} [Proceedings of the V.\,V.~Vinogradov Russian Language 


Institute] 6:300--310.


\bibitem{22-gz-1} %15


\Aue{Zaliznyak, Anna~A., D.\,V.~Sitchinava, S.~Loiseau, M.~Kruzhkov, and 


I.\,M.~Zatsman.} 2013. Database of equivalent verbal forms in a~Russian-French 


multivariant parallel corpus. \textit{Conference (International) on Artificial 


Intelligence Proceedings}.


Las Vegas, NV: CSREA Press. 1:101--107.


\bibitem{20-gz-1} %16


\Aue{Kruzhkov, M.\,G., N.\,V.~Buntman, E.\,Ju.~Loshchilova, D.\,V.~Sitchinava, 


Anna A.~Zalizniak, and I.\,M.~Zatsman.} 2014. A~database of Russian verbal forms 


and their French translation equivalents. \textit{Komp'yuternaya lingvistika 


i~intellektual'nye tekhnologii: po mat-lam ezhegodnoy Mezhdunar. konf. ``Dialog''} 


[Computational Linguistics and Intellectual Technologies. Papers from the Annual 


Conference (International) ``Dialogue'']. Moscow: RGGU. 13(20):275--287.


\bibitem{21-gz-1} %17


\Aue{Buntman, N.\,V., Anna~A.~Zaliznyak, I.\,M.~Zatsman, M.\,G.~Kruzhkov, 


E.\,Yu.~Loshchilova, and D.\,V.~Sichinava.} 2014. Informatsionnye tekhnologii 


korpusnykh issledovaniy: printsipy postroeniya krosslingvisticheskikh baz dannykh 


[Information technologies for corpus studies: Underpinnings for cross-linguistic 


database creation]. \textit{Informatika i~ee Primeneniya~--- Inform. Appl.}  


8(2):98--110.





\bibitem{17-gz-1} %18


\Aue{Popkova, N.\,A., O.\,Yu.~Inkova, I.\,M.~Zatsman, and M.\,G.~Kruzhkov.} 


2015. Metodika postroeniya monoekvivalentsiy v~nadkorpusnoy baze dannykh 


konnektorov [Methodology of constructing monoequivalences in the supracorpora 


database of connectors]. \textit{Tr. 2-y molodezhnoy nauchn. konf. ``Zadachi sovremennoy 


informatiki''} [2nd Scientific Conference ``Problems of Modern Informatics'' 


Proceedings]. Moscow: FRC CSC RAS. 143--153.


\bibitem{18-gz-1} %19


\Aue{Zaliznyak, Anna~A., I.\,M.~Zatsman, O.\,Yu.~Inkova, and M.\,G.~Kruzhkov.} 


2015. Nadkorpusnye bazy dannykh kak lingvisticheskiy resurs [Supracorpora 


databases as linguistic resource]. 


\textit{Conference (International) ``Corpus Linguistics-2015'' 


Proceedings}. St.\ Petersburg: St. Petersburg State University. 211--218.


\bibitem{19-gz-1} %20


\Aue{Kruzhkov, M.\,G.} 2015. Informatsionnye resursy kontrastivnykh 


lingvisticheskikh issledovaniy: elektronnye korpusa tekstov [Information resources 


for contrastive studies: Electronic text corpora]. \textit{Sistemy i Sredstva 


Informatiki~--- Systems and Means of Informatics} 25(2):140--159.





\bibitem{15-gz-1} %21


\Aue{Inkova, O.\,Yu., and M.\,G.~Kruzhkov.} 2016. Nadkorpusnye  


russko-frantsuzskie basy dannykh glagol'nykh form konnektorov [Supracorpora 


databases of Russian and French verbal forms and connectors]. \textit{Lingue slave 


a~confronto}. Eds. O.~Inkova and A.~Trovesi. 


Bergamo: Bergamo University Press. 365--392.


\bibitem{16-gz-1} %22


\Aue{Zaliznyak, Anna A., and M.\,G.~Kruzhkov.} 2016. Baza dannykh bezlichnykh 


glagol'nykh konstruktsiy russkogo yazyka [Database of Russian impersonal verbal 


constructions]. \textit{Informatika i~ee Primeneniya~--- Inform. Appl.}  


10(4):132--141.











\bibitem{23-gz-1}


\Aue{Inkova-Manzotti, O.\,Yu.} 2001. \textit{Konnektory protivopostavleniya vo 


frantsuzskom i~russkom yazykakh (sopostavitel'noe issledovanie)} [Connectives of 


opposition in French and Russian (comparative research)]. Moscow: Informelektro. 


434~p.


\bibitem{24-gz-1}


\Aue{Zatsman, I.\,M.} 2019. Tselenapravlennoe razvitie sistem lingvisticheskikh 


znaniy: vyyavlenie i zapolnenie lakun [Goal-oriented development of linguistic 


knowledge systems: Identifying and filling lacunae]. \textit{Informatika i~ee 


Primeneniya~--- Inform. Appl.} 13(1):91--98.


\end{thebibliography}


} }








\vspace*{-3pt}





\hfill{\small\textit{Received February 15, 2019}}








\Contr





\noindent


\textbf{Goncharov Alexander A.} (b.\ 1994)~--- junior scientist, Institute of 


Informatics Problems, Federal Research Center ``Computer Science and Control'' 


of the Russian Academy of Sciences, 44-2~Vavilov Str., Moscow 119333, Russian 


Federation; \mbox{a.gonch48@gmail.com}





\vspace*{3pt}





\noindent


\textbf{Zatsman Igor M.} (b.\ 1952)~--- Doctor of Science in technology, Head of 


Department, Institute of Informatics Problems, Federal Research Center 


``Computer Science and Control'' of the Russian Academy of Sciences,  


44-2~Vavilov Str., Moscow 119333, Russian Federation; \mbox{izatsman@yandex.ru}





\label{end\stat}








\renewcommand{\bibname}{\protect\rm Литература}


